\clearpage\fancyhead[L]{\small NEW ELEMENTS OF PALLAS.}\SourceLeaf{325}

\begin{center}\textbf{NEW ELEMENTS OF PALLAS.}\par\vspace{4pt}\hrule\vspace{4pt}\fontsize{10}{12}\selectfont Monatliche Correspondenz. Vol. XXII. pp. 400\,..\,403. October 1810.\par\vspace{4pt}\hrule\vspace{4pt}\end{center}

\par\indent Of the new planets, \emph{Pallas} will be the first to become visible and will be quite bright because it will be so close to Earth. In our northern regions its brightness might be reduced only by its large southern declination. We therefore hasten to share with our readers the ephemeris of its motion recently calculated by Professor \textsc{Gauss}. At the beginning of November it will probably already be possible to observe the planet on the meridian.

The elements on which the following ephemeris is based were derived from the oppositions of 1805, 7, 8 and 9. They are given below.

\Needspace{14\baselineskip}\begin{center}\emph{Göttingen meridian.}\par\medskip\fontsize{10}{12}\selectfont\setlength{\tabcolsep}{4pt}\renewcommand{\arraystretch}{1.15}\begin{tabular}{llr}
Mean longitude 1803 & .\;.\;. & $221^\circ\;23'\;24''6$\\
mean daily motion &  & $770''9365$\\
1803 perihelion & .\;.\;.\;. & $120^\circ\;58'\;4''8$\\
--- $\AscendingNode$ & .\;.\;.\;.\;.\;. & $172\;27\;52.4$\\
Inclination & .\;.\;,\;.\;.\;.\;. & $34^\circ\;36'\;49''4$\\
Eccentricity & .\;.\;.\;.\;. & $0.2446335$\\
Log. semimajor axis & .\;.\;.\;. & $0.4420473$\\
\end{tabular}\end{center}

\par\indent [The motion of Pallas in 1810 and 1811 is given for midnight in Göttingen at four-day intervals, R.A. and decl. in minutes log. dist. from Earth to four decimal places]
